% Integración por delta de la adenda de realizaciones físicas.
% Residencia principal: transferencia física de las tres álgebras y lector temporal.

\chapter{Quadratic regimes and temporal readout}
\label{chap:regimenes-cuadraticos-tiempo}
\index{quadratic algebra!three regimes}
\index{time crystal!finite HMT realization}
\index{time!TRIT reader}

\section{Quadratic algebras of the TRIT and curvature signatures}
\label{sec:regimenes-cuadraticos-recibidos}

The primary construction and proof of the three regimes reside in
\cref{chap:trit-algebras}. For
\(\tau\in\TRIT\), with balanced image
\(\bar\tau\in\{-1,0,+1\}\), one has
\[
 \mathbb A_\tau=\RR[X]/(X^2+\bar\tau),
\]
and, if \(J\) is the class of \(X\),
\[
\begin{array}{ccl}
 \tau=+1&:&J^2=-I,\quad\text{elliptic regime},\\
 \tau=0&:&J^2=0,\quad\text{parabolic regime},\\
 \tau=-1&:&J^2=+I,\quad\text{hyperbolic regime}.
\end{array}
\]
For all three values of \(\tau\), the quotient
\(\mathbb A_\tau=\mathbb R[X]/(X^2+\bar\tau)\) satisfies precisely the three
quadratic identities above. In particular, the internal square root of
\(-1\) is realized as an operator on the elliptic sheet; HMT need neither reject nor
suppress the complex extension. The parabolic algebra contains a nilpotent,
and the hyperbolic algebra decomposes by means of the two idempotents
\((I\pm J)/2\).

These identities classify algebras and transports. They do not, by themselves,
select a spacetime metric or prove that the physical universe occupies one
of the three sheets.

\section{Parametrized metric family of the temporal reader}
\label{sec:ansatz-metrico-tres-curvaturas}

Write
\[
 \mathbb A_s:=\RR[X]/(X^2+s),
 \qquad s\in\{+1,0,-1\},
\]
for the same family with the real sign as index.
To make the possible interface precise, fix a scale
\(\rho_{\rm curv}>0\), of dimension \(L^{-1}\), and
\(\theta\in\RR/2\pi\ZZ\). For \(s\in\{+1,0,-1\}\), let
\[
 g_{s,\rho}
 =\dd r^2+S_{s,\rho}(r)^2\,\dd\theta^2,
\]
\[
 S_{s,\rho}(r)=
 \begin{cases}
  \rho_{\rm curv}^{-1}\sin(\rho_{\rm curv}r),&s=+1,\\
  r,&s=0,\\
  \rho_{\rm curv}^{-1}\sinh(\rho_{\rm curv}r),&s=-1.
 \end{cases}
\]
The radial domain is \(r\geq0\) for \(s=0,-1\); for \(s=+1\), the chart
polar cubre \(0\leq r\leq\pi/\rho_{\rm curv}\), with the polar regularity
usual in the extremos. In the interior of each chart,
\[
 K_{\rm G}
 =-\frac{S_{s,\rho}''}{S_{s,\rho}}
 =s\,\rho_{\rm curv}^{\,2}.
\]
The last equality is an exact computation of the model once the realization is fixed. The arrow
\[
 \mathbb A_s\dashrightarrow(M_s,g_{s,\rho})
\]
is, however, a \emph{ansatz} of the \textsc{physical interface}; it is not a
canonical consequence of the relation \(J^2=-sI\).

If \(\ell_0\in\mathcal L_L^+\) is a length and we declare
\[
 \rho_{\rm curv}:=\frac{\Delta_4}{\ell_0},
\]
then
\[
 K_{\rm G}
 =s\,\frac{\Delta_4^2}{\ell_0^2}.
\]
This specialization is
\textsc{demonstrated with explicit starting structure}: the conclusion is
obtained once \(s,\Delta_4,\ell_0\) have been declared, but it does not select
\(\ell_0\), does not
deduce the observed cosmological curvature, and does not identify the ratios
\(\mathcal G_{\rm tor}/\mathcal I_{\rm tor}\) or
\(\mathcal G_{\rm av}/\mathcal I_{\rm av}\) with \(K_{\rm G}\).

\section{Temporal TRIT reader on ({-1,0,+1}): opening, threshold, and return with memory}
\label{sec:lector-temporal-trit}

Once an ordered parametrization of the transitions has been chosen, the active
temporal reader is
\[
\begin{aligned}
 +1&\longmapsto
 \text{return, retained memory and past},\\
 0&\longmapsto
 \text{evaluation threshold and present},\\
 -1&\longmapsto
 \text{bidirectional opening and future}.
\end{aligned}
\]
The reader is \textsc{proved with an explicit starting structure}: its
conclusions follow once the TRIT signature, the order of the
trajectory, and the transported memory have been specified. It does not turn the set
\(\{-1,0,+1\}\) into an equation of motion. The future designates the space
of compatible continuations; survival may select a branch
or a multisection without erasing the genealogy of the others.

The frequency
\[
 \omega_0=\frac{2\pi}{108\,t_0}
\]
assigns a type to the full-return clock through the decision entropy
\(I_{\mathrm{TRIT}}=\log 3\) and the Boltzmann transducer.
It does not, in isolation, establish a stable subharmonic phase. It does, however, enable construction of
the finite reference Hamiltonian below; selection of a robust physical
Hamiltonian requires additional conditions.

\section{Unitary dynamics on 108 states: Hamiltonian, propagator, and period}
\label{sec:hamiltoniano-reloj-108}

Let
\[
 \mathcal H_{\rm clk}:=\mathbb C^{108},
 \qquad
 \{|j\rangle:j\in\mathbb Z/108\mathbb Z\}
\]
with its canonical Hermitian inner product. Fix
\(\zeta:=\exp(2\pi\mathrm i/108)\), the unitary shift
\[
 U_{\rm clk}|j\rangle=|j+1\rangle
\]
and the Fourier basis
\[
 |\widetilde k\rangle
 :=\frac1{\sqrt{108}}\sum_{j=0}^{107}\zeta^{jk}|j\rangle,
 \qquad 0\leq k\leq107.
\]
With \(t_0\in\mathcal L_T^+\) and
\(\hbar_{\rm int}\in\mathcal L_S^+\), define
\[
 \boxed{
 H_{\rm clk}
 :=\hbar_{\rm int}\omega_0
 \sum_{k=0}^{107}k\,
 |\widetilde k\rangle\langle\widetilde k|,
 \qquad
 \omega_0=\frac{2\pi}{108\,t_0}.}
 \tag{CLK1}\label{eq:hamiltoniano-reloj-108}
\]
The operator has energy type
\([H_{\rm clk}]=\mathcal L_S\mathcal L_T^{-1}=\mathcal L_E\).

\begin{theorem}[Exact unitary realization of the \(108\) calendar]
\label{thm:reloj-unitario-108}
The operator \(H_{\rm clk}\) is self-adjoint and its one-step propagator
satisfies
\[
 U_{\rm step}
 :=\exp\!\left(-\frac{\mathrm i}{\hbar_{\rm int}}
 H_{\rm clk}t_0\right)
 =U_{\rm clk},
 \qquad
 U_{\rm step}^{108}=I.
 \tag{CLK2}\label{eq:propagador-reloj-108}
\]
If
\[
 Z_{\rm clk}|j\rangle=\zeta^j|j\rangle,
\]
then, for any basis state \(|j_0\rangle\),
\[
 \langle Z_{\rm clk}\rangle_{n}
 :=\langle j_0|U_{\rm step}^{-n}Z_{\rm clk}
 U_{\rm step}^n|j_0\rangle
 =\zeta^{j_0+n}.
 \tag{CLK3}\label{eq:respuesta-reloj-108}
\]
The sequence has primitive period \(108\).
\end{theorem}

\begin{proof}
Expression \eqref{eq:hamiltoniano-reloj-108} is a spectral
decomposition with real eigenvalues, hence \(H_{\rm clk}=H_{\rm clk}^\dagger\).
Moreover,
\[
 U_{\rm clk}|\widetilde k\rangle
 =\zeta^{-k}|\widetilde k\rangle,
\]
whereas the exponential of
\(\hbar_{\rm int}\omega_0 k\) over \(t_0\) gives
\(\exp(-2\pi\mathrm ik/108)=\zeta^{-k}\). The two operators agree on
an orthonormal basis. The remaining identities follow from
\(U_{\rm step}^n|j_0\rangle=|j_0+n\rangle\) and the primitivity of
\(\zeta\).
\end{proof}

The corresponding variational formulation, for a curve
\(\psi\in C^1([t_a,t_b],\mathcal H_{\rm clk})\), unconstrained during the
variation, is
\[
 \boxed{
 L_{\rm clk}(\psi,\dot\psi)
 =\frac{\mathrm i\hbar_{\rm int}}2
 \bigl(
 \langle\psi|\dot\psi\rangle
 -\langle\dot\psi|\psi\rangle
 \bigr)
 -\langle\psi|H_{\rm clk}|\psi\rangle.}
 \tag{CLK4}\label{eq:lagrangiano-reloj-108}
\]
The action functional is
\[
 \boxed{
 S_{\rm clk}[\psi]
 :=\int_{t_a}^{t_b}
 L_{\rm clk}(\psi(t),\dot\psi(t))\,\mathrm dt,}
 \tag{CLK5}\label{eq:accion-reloj-108}
\]
with \(\delta\psi(t_a)=\delta\psi(t_b)=0\).
\index{Schrödinger!finite-clock equation}
Varying \(\psi\) and \(\bar\psi\) independently yields the
Schrödinger equation for the finite clock
\[
 \boxed{\mathrm i\hbar_{\rm int}\dot\psi=H_{\rm clk}\psi}
 \tag{CLK6}\label{eq:schrodinger-reloj-108}
\]
together with its adjoint equation \cite{SchrodingerNobel}. Its propagator sampled
at \(t=nt_0\) is \(\psi_n=U_{\rm step}^{n}\psi_0\). Since \(H_{\rm clk}\) is
self-adjoint, the norm is conserved; normalized initial data remain
normalized. The result is exact with respect to the specified
\(H_{\rm clk},t_0,\hbar_{\rm int}\); it neither selects a physical Hamiltonian
nor introduces an experimental frequency.

The \cref{thm:reloj-unitario-108} proves the unitary dynamics of the calendar relative to the declared complex completion, \(t_0\) and \(\hbar_{\rm int}\). The five dynamical conditions of the following section receive a constructive realization in \cref{sec:hmtf2-cristal-temporal}: there the drive, the order observable, the primitive subharmonic of period \(108T_{\rm d}\), and finite spectral stability are fixed. The macroscopic material realization constitutes a distinct external recognition and not a lacuna in the finite result.

\section{Dynamic realization conditions of a time crystal: Hamiltonian, subharmonic response and stability}
\label{sec:programa-cristal-tiempo}

A driven realization must provide, at a minimum:
\begin{enumerate}
 \item a physical state space and a self-adjoint Hamiltonian
 \(H(t)\), or an equivalent driven protocol, with
 \(H(t+T_{\rm d})=H(t)\);
 \item the propagator of Floquet
 \[
  U_F=\mathcal T\exp\!\left(
  -\frac{\mathrm i}{\hbar_{\rm int}}
  \int_0^{T_{\rm d}}H(t)\,\dd t\right);
 \]
 \item an observable of order \(\mathcal O\) and a family of states for
 which there exists a primitive subharmonic response of order \(m>1\):
 \[
  \langle\mathcal O((n+m)T_{\rm d})\rangle
  =\langle\mathcal O(nT_{\rm d})\rangle
 \]
 in the asymptotic regime, without the same periodicity for a proper divisor
 of \(m\);
 \item stability of \(m\), of the phase, and of the amplitude under an open
 class of admissible perturbations;
 \item persistence that is not a finite-size transient, together with a
 declared mechanism of localization, prethermalization, or dynamical
 protection.
\end{enumerate}
The calendar \(27\to54\to108\) provides the internal hierarchy of periods. The
finite realization of \cref{sec:hmtf2-cristal-temporal} selects subharmonic order
\(108\) in its declared domain, without replacing the five preceding
constructions or by itself imposing a macroscopic material phase.

\begin{criterion}[HMT--MD time crystal]
\label{crit:cristal-tiempo-hmt-md}
The identification is refuted if no well-typed driven Hamiltonian or protocol
exists; if the alleged subharmonicity disappears under
arbitrarily small perturbations compatible with the model; if the
observed period is only the period imposed by the drive; if the
oscillation is a finite transient; or if the temporal reader assigns equivalent
histories to incompatible physical phases.
\end{criterion}

\paragraph{Mathematical and physical domain of the temporal reader.}
The three algebras and their exponentials are proved in the algebraic domain; the temporal reader is deduced from the initial signature, order, and memory. The finite HMT time crystal satisfies the \cref{crit:cristal-tiempo-hmt-md} through the construction of \cref{sec:hmtf2-cristal-temporal}. Extension to a macroscopic material phase additionally requires specification of the physical support and belongs to the subsequent experimental comparison.
